\documentclass[11pt,letterpaper]{article}
\usepackage[T1]{fontenc}
\usepackage[utf8]{inputenc}
\usepackage{lmodern}
\usepackage[margin=1in,headheight=14pt]{geometry}
\usepackage{amsmath,amssymb,amsthm,mathtools}
\usepackage{microtype}
\usepackage{booktabs,tabularx,longtable,array}
\usepackage{enumitem}
\usepackage{needspace}
\usepackage{fancyhdr}
\usepackage{xcolor}
\usepackage{hyperref}
\hypersetup{colorlinks=true,linkcolor=blue!45!black,citecolor=blue!45!black,
 urlcolor=blue!45!black,pdftitle={Skew partitions, continued fractions, and connected components},
 pdfauthor={ChatGPT},pdfsubject={A complete proof of the generating function recorded as conjectural in OEIS A225114}}
\usepackage{bookmark}
\numberwithin{equation}{section}
\newtheorem{theorem}{Theorem}[section]
\newtheorem{proposition}[theorem]{Proposition}
\newtheorem{lemma}[theorem]{Lemma}
\newtheorem{corollary}[theorem]{Corollary}
\theoremstyle{definition}
\newtheorem{definition}[theorem]{Definition}
\newtheorem{example}[theorem]{Example}
\newtheorem{remark}[theorem]{Remark}
\newcommand{\Z}{\mathbb Z}
\newcommand{\N}{\mathbb N}
\newcommand{\C}{\mathbb C}
\newcommand{\Aall}{\mathcal A}
\newcommand{\Pconn}{P}
\newcommand{\comp}{\operatorname{comp}}
\newcommand{\E}{\mathbb E}
\newcommand{\Var}{\operatorname{Var}}
\newcommand{\opnorm}[1]{\left\lVert#1\right\rVert}
\newcommand{\coeff}[2]{\left[#1\right]#2}
\newcolumntype{Y}{>{\raggedright\arraybackslash}X}
\pagestyle{fancy}
\fancyhf{}
\fancyhead[L]{\small Skew partitions and continued fractions}
\fancyhead[R]{\small OEIS A225114}
\fancyfoot[C]{\thepage}
\setlength{\parskip}{3pt}
\setlength{\emergencystretch}{2em}
\setlist{nosep}

\title{\vspace{-1.2em}\textbf{Skew partitions, continued fractions,\\
 and connected components}\\[0.65em]
 \large A complete proof of the generating function\\
 recorded as conjectural in OEIS A225114}
\author{A research note prepared with ChatGPT}
\date{20 September 2026}

\begin{document}
\maketitle
\thispagestyle{plain}

\begin{abstract}
Let $a_n$ count skew partition diagrams with $n$ cells and no empty rows or
columns, including one empty diagram. We prove that
\[
 \sum_{n\ge0}a_nq^n=\frac{1}{2-C(q)},\qquad
 C(q)=\cfrac{1}{1-\cfrac{q}{1-\cfrac{q}{1-\cfrac{q^2}{1-\cfrac{q^2}{1-\ddots}}}}},
\]
where every power $q^j$ occurs twice. The proof gives explicit bijections from
row-overlap data to diagrams, from connected diagrams to pairs of boundary
paths, and from those pairs to weighted two-colored Motzkin excursions.
It also establishes refinements by area, rows, columns, and connected
components; a closed formula at fixed dimensions; an independent ratio of
$q$-series; a sharp square-law truncation bound; and exponential asymptotics
with a rationally certified dominant-singularity interval.

The target was selected by one recorded random draw from a frozen shortlist.
Although OEIS still describes its generating function as conjectural, Dima
Pasechnik published a proof sketch in September 2024. Accordingly, this is a
self-contained proof and elaboration, not a claim to the first solution of
an unresolved conjecture. Exact independent checks through area $1000$ and
all source code accompany the article.
\end{abstract}

\newpage
\begingroup
\small
\setlength{\parskip}{0pt}
\tableofcontents
\endgroup
\newpage

\section{Selection, provenance, and the main result}
\label{sec:selection}

\subsection{The frozen candidate list}
The search was restricted to explicit conjectured ordinary generating
functions for natural combinatorial or sequence-theoretic objects. Prime
number and congruence questions were excluded. The following four candidates
were fixed before a random index was generated. This is a screened shortlist,
not a claim to survey every eligible OEIS conjecture.

\begin{table}[htbp]
\centering\small
\begin{tabularx}{\textwidth}{@{}c l Y l@{}}
\toprule
Index & Entry & Subject & Conjecture attribution\\
\midrule
0 & A225114 & Skew partitions without empty rows or columns & Kurkov, 2024\\
1 & A244475 & Fifth-largest distinct value in a row of Stern's triangle & Heinz, 2022\\
2 & A381190 & Connected minimal dominating sets of trapezohedral graphs & Arndt, 2026\\
3 & A391632 & Mutual-visibility sets of Andr\'asfai graphs & Howroyd, 2026\\
\bottomrule
\end{tabularx}
\caption{The ordered list used for the draw. The entries and their explicit
formulas were checked on OEIS \cite{oeis225114,oeis244475,oeis381190,oeis391632}.
The three unselected formulas are reproduced in Appendix~\ref{app:candidates}.}
\end{table}

The recorded method was \texttt{random.Random(seed).randrange(4)}, with a
128-bit seed obtained from \texttt{secrets.randbits(128)}. The seed was
\[
 \texttt{203225573310603175027560032786773102290}.
\]
The result was the zero-based index $0$, hence \textbf{A225114}. There was one
draw and no reroll. The files \texttt{candidates.json} and
\texttt{selection.json} preserve the list and result; the verification
script checks the list's digest and replays the draw. This records the actual
selection procedure, not merely a seed chosen after selecting a target.

\subsection{What is being proved, and what was already known}
The selected entry counts skew partitions whose diagrams have no empty rows
or columns. It attributes the conjectured continued fraction to Mikhail
Kurkov, 3 September 2024 \cite{oeis225114}. These diagrams connect a familiar
partition-theoretic construction with lattice-path and polyomino enumeration;
they are not an artificially defined numerical sequence.

A subsequent prior-art check found Pasechnik's MathOverflow answer of
10 September 2024 \cite{pasechnik}. That answer sketches the generating-function
identification and the ordered-sequence construction. The connected-diagram
continued fraction is older: it appears in the record for A006958 and in
work on skew Ferrers diagrams \cite{oeis006958,bala}. The present article
therefore retains the randomly selected target, proves its formula in full,
and supplies additional derivations and consequences. It does not assert
novelty of the connected continued fraction, the sequence construction, or
the underlying enumerative identities.

\Needspace{14\baselineskip}
\begin{theorem}[The selected generating function]
\label{thm:main}
Let $a_n$ be the number of normalized skew partition diagrams of area $n$
with no empty rows or columns, with $a_0=1$. Define
\begin{equation}
 C(q)=\cfrac{1}{1-\cfrac{q}{1-\cfrac{q}{1-\cfrac{q^2}{1-\cfrac{q^2}{1-
 \cfrac{q^3}{1-\cfrac{q^3}{1-\ddots}}}}}}}.
 \label{eq:mainC}
\end{equation}
The continued fraction is well defined coefficientwise in $\Z[[q]]$, and
\begin{equation}
 A(q):=\sum_{n\ge0}a_nq^n=\frac{1}{2-C(q)}.
 \label{eq:mainA}
\end{equation}
\end{theorem}

The essential decomposition is
\begin{equation}
 \boxed{\quad A(q)=\frac{1}{1-P(q)},\qquad C(q)=1+P(q),\quad}
 \label{eq:proofplan}
\end{equation}
where $P(q)$ counts nonempty \emph{edge-connected} diagrams. Neither equality
will be inferred from a finite list of coefficients. Sections~\ref{sec:rows}
and~\ref{sec:paths} prove the geometric constructions, and
Section~\ref{sec:contraction} proves the continued-fraction identity.

\section{Row lengths, overlaps, and connected components}
\label{sec:rows}

\subsection{Normalization and the exact counting convention}

\begin{definition}
A nonempty normalized skew partition diagram is specified by partitions
$\lambda=(\lambda_1,\ldots,\lambda_h)$ and
$\mu=(\mu_1,\ldots,\mu_h)$, where $\mu$ is padded with zeros,
\[
 \lambda_i>\mu_i\ge0,\qquad \mu_h=0.
\]
Its cells in row $i$, with rows read from top to bottom, have column indices
$\mu_i+1,\ldots,\lambda_i$. Its area, number of rows, and number of columns
are denoted $n,h,w$, respectively. Every column from $1$ to $w=\lambda_1$
must contain a cell. The empty diagram is included separately.
\end{definition}

The normalization removes horizontal translation. Empty rows are absent by
$\lambda_i>\mu_i$. Diagrams are not identified under reflection,
transposition, or rotation. We always use $\lambda$ for the outer partition
and $\mu$ for the inner partition.

Put
\begin{equation}
 \ell_i=\lambda_i-\mu_i>0,\qquad
 o_i=\lambda_{i+1}-\mu_i\quad(1\le i<h).
 \label{eq:lengthoverlap}
\end{equation}
The number $o_i$, when nonnegative, is the number of columns shared by rows
$i$ and $i+1$. In particular, $o_i>0$ means that those rows share a horizontal
cell edge. The case $o_i=0$ is corner contact, not edge connectivity.

\begin{lemma}[The no-empty-column condition]
\label{lem:noempty}
A diagram as above has no empty columns if and only if $o_i\ge0$ for every
$1\le i<h$.
\end{lemma}
\begin{proof}
If $\lambda_{i+1}<\mu_i$, every column with index strictly larger than
$\lambda_{i+1}$ and at most $\mu_i$ is empty. Rows above the cut start to its
right, since their inner-partition parts are at least $\mu_i$; rows below
the cut end to its left, since their outer-partition parts are at most
$\lambda_{i+1}$.

Conversely, suppose $o_i\ge0$. The two consecutive row intervals overlap or
are adjacent, because
$\mu_i+1\le\lambda_{i+1}+1$. Their union therefore contains no missing
integer column. Inducting upwards from the bottom row, which starts in
column $1$, shows that the union of all row intervals is precisely
$\{1,\ldots,\lambda_1\}$. Thus no column is empty.
\end{proof}

\begin{proposition}[Independent overlap coordinates]
\label{prop:rowcoords}
Fix positive row lengths $(\ell_1,\ldots,\ell_h)$. Normalized diagrams with
those row lengths and no empty columns are in bijection with choices
\begin{equation}
 0\le o_i\le\min(\ell_i,\ell_{i+1})\qquad(1\le i<h).
 \label{eq:overlapbounds}
\end{equation}
Every choice in these ranges is possible, independently of the others.
\end{proposition}
\begin{proof}
For an existing diagram, the two partition inequalities give
\[
 o_i\le\lambda_i-\mu_i=\ell_i,
 \qquad o_i\le\lambda_{i+1}-\mu_{i+1}=\ell_{i+1},
\]
and Lemma~\ref{lem:noempty} gives the lower bound.

Conversely, begin with $\mu_h=0$ and $\lambda_h=\ell_h$. For
$i=h-1,h-2,\ldots,1$, define
\begin{equation}
 \mu_i=\lambda_{i+1}-o_i,\qquad
 \lambda_i=\mu_i+\ell_i.
 \label{eq:reconstruction}
\end{equation}
The upper bounds in \eqref{eq:overlapbounds} imply
\[
 \mu_i-\mu_{i+1}=\ell_{i+1}-o_i\ge0,
 \qquad
 \lambda_i-\lambda_{i+1}=\ell_i-o_i\ge0.
\]
Thus $\lambda$ and $\mu$ are partitions, all rows are nonempty, and
Lemma~\ref{lem:noempty} applies. Equations~\eqref{eq:reconstruction} also
prove uniqueness.
\end{proof}

\begin{corollary}[Finite general-term sums]
\label{cor:compositions}
For $n\ge1$,
\begin{align}
 a_n&=\sum_{h=1}^{n}\ \sum_{\substack{\ell_1+\cdots+\ell_h=n\\\ell_i\ge1}}
       \prod_{i=1}^{h-1}\bigl(\min(\ell_i,\ell_{i+1})+1\bigr),
       \label{eq:allcomposition}\\
 p_n&=\sum_{h=1}^{n}\ \sum_{\substack{\ell_1+\cdots+\ell_h=n\\\ell_i\ge1}}
       \prod_{i=1}^{h-1}\min(\ell_i,\ell_{i+1}),
       \label{eq:connectedcomposition}
\end{align}
where $p_n$ counts edge-connected diagrams and an empty product equals $1$.
Moreover, if $a_n(u)$ marks the number of connected components, then
\begin{equation}
 a_n(u)=\sum_{h=1}^{n}\ \sum_{\substack{\ell_1+\cdots+\ell_h=n\\\ell_i\ge1}}
 u\prod_{i=1}^{h-1}\bigl(\min(\ell_i,\ell_{i+1})+u\bigr).
 \label{eq:componentcomposition}
\end{equation}
\end{corollary}
\begin{proof}
Each positive overlap joins its two adjacent rows. Each zero overlap
separates two components, and every component consists of consecutive rows.
Thus the number of components is one plus the number of zero overlaps.
At each boundary there are $\min(\ell_i,\ell_{i+1})$ positive choices and
one zero choice; weighting the latter by $u$, and the first component by
$u$, gives all three formulas.
\end{proof}

The first two sums also explain the composition-based program already
recorded in the OEIS entry \cite{oeis225114}. They are useful independent
checks, but the desired continued fraction is not yet evident from them.

\begin{example}
For row lengths $(2,2)$ the overlaps $0,1,2$ give, respectively,
\[
 (\lambda,\mu)=((4,2),(2,0)),\quad
 ((3,2),(1,0)),\quad ((2,2),(0,0)).
\]
The first diagram has two components. The other two are connected. Their
widths are $4,3,2$, illustrating the general identity
\begin{equation}
 w=\lambda_1=n-\sum_{i=1}^{h-1}o_i.
 \label{eq:widthoverlap}
\end{equation}
This follows by telescoping
$\lambda_i=\lambda_{i+1}+\ell_i-o_i$.
\end{example}

\subsection{The ordered-sequence construction}
Let
\begin{align}
 P(q;s,t)&=\sum_{\substack{D\text{ nonempty}\\D\text{ connected}}}
               q^{n(D)}s^{h(D)}t^{w(D)},\label{eq:connectedrefined}\\
 \Aall(q;s,t,u)&=\sum_D q^{n(D)}s^{h(D)}t^{w(D)}u^{k(D)},
 \label{eq:allrefined}
\end{align}
where the empty diagram has weight $1$ and $k(D)$ is the number of
edge-connected components. Write $P(q)=P(q;1,1)$, so $p_0=0$.

\begin{proposition}[Unique concatenation into components]
\label{prop:sequence}
There is a weight-preserving bijection between all normalized diagrams and
finite ordered sequences of nonempty connected normalized diagrams.
Consequently,
\begin{equation}
 \Aall(q;s,t,u)=\frac{1}{1-uP(q;s,t)}.
 \label{eq:sequence}
\end{equation}
\end{proposition}
\begin{proof}
Cut a diagram between rows $i$ and $i+1$ whenever $o_i=0$, then translate
each resulting block horizontally to normalize it. Positive overlaps
remain within every block, so these blocks are exactly its components.
Their order is the order of the row blocks, from top to bottom.

Conversely, concatenate any ordered collection of connected blocks. Retain
their internal row lengths and positive overlaps, and insert a zero overlap
between each consecutive pair. Proposition~\ref{prop:rowcoords} reconstructs
a unique normalized diagram. Geometrically, each upper block is placed
northeast of the next block, with their bounding boxes touching at one
corner. These two constructions are inverse.

Areas and row counts add. Widths add as well: the inserted overlaps are
zero, so \eqref{eq:widthoverlap} is additive over blocks. A sequence of $k$
blocks therefore has generating function $u^kP^k$. Summing over $k\ge0$
gives \eqref{eq:sequence}; this is legitimate since $P$ has zero constant
term in $q$.
\end{proof}

\section{Connected diagrams as weighted Motzkin excursions}
\label{sec:paths}

\subsection{Two boundary paths}
Place the diagram in the Cartesian plane, with the bottom row at height
zero. A connected diagram of height $h$ and width $w$ is bounded by two
north-east lattice paths $U$ and $L$ from $(0,0)$ to $(w,h)$. Here $U$ is
the upper or northwest boundary, and $L$ is the lower or southeast boundary.
The first steps are, respectively, north and east; the last steps are east
and north. The two paths have no common point except their endpoints.

For completeness, the partition inequalities make both boundary paths
north-east monotone. A meeting at an intermediate horizontal row boundary
would say $\lambda_{i+1}=\mu_i$, namely a zero overlap. A shared horizontal
boundary segment would force an empty row. Thus positive overlaps and
nonempty rows imply that the paths meet only at their endpoints.
Conversely, the cells between such a pair of paths have contiguous rows;
their left and right endpoints are monotone as rows are read upwards.
They therefore form a skew partition diagram, with positive overlaps and
no empty rows or columns. This proves the boundary-path bijection.

Both paths have length $L_0=h+w$. At integer time $j$, meaning on the line
$x+y=j$, let
\begin{equation}
 d_j=y_U(j)-y_L(j).
 \label{eq:gaps}
\end{equation}
Then $d_0=d_{L_0}=0$, and $d_j\ge1$ for $1\le j<L_0$.

\begin{lemma}[Area from synchronous separation]
\label{lem:area}
The area of the diagram is
\begin{equation}
 n=\sum_{j=1}^{L_0-1}d_j.
 \label{eq:areasum}
\end{equation}
\end{lemma}
\begin{proof}
Apply the linear map $(x,y)\mapsto(x+y,y)$, whose determinant is $1$.
The area between the paths becomes the integral of their vertical gap as
a function of the new first coordinate. On every interval $[j,j+1]$ this
gap is linear, and its integral is $(d_j+d_{j+1})/2$. Summation gives
\eqref{eq:areasum}, because the endpoint gaps are zero. The original area
is the number of unit cells.
\end{proof}

\subsection{The four step pairs and their weights}
The pairs of simultaneous boundary steps have the following effects:
\begin{center}
\begin{tabular}{@{}c c c c@{}}
\toprule
Upper step & Lower step & Change of gap & Dimension weight\\
\midrule
North & East & $+1$ & $s$\\
East & North & $-1$ & $t$\\
North & North & $0$, first color & $s$\\
East & East & $0$, second color & $t$\\
\bottomrule
\end{tabular}
\end{center}
Dimension weights can be assigned using the upper path alone: its number
of north steps is $h$, and its number of east steps is $w$.

Delete the forced first up-step and last down-step in the gap process,
then lower all remaining gaps by $1$. The result is a two-colored Motzkin
excursion: a path starting and ending at height zero, staying nonnegative,
and using up-steps, down-steps, and two kinds of horizontal steps.
Conversely, the table recovers both boundary words from any such excursion,
after the two forced endpoint pairs are restored. There is no ambiguity in
either direction.

Give the two removed endpoint steps together the factor $stq$. For an
internal step landing at height $b$ in the Motzkin path, give an additional
factor $q^{b+1}$, together with its dimension weight from the table.
Lemma~\ref{lem:area} proves that the full product is exactly $q^ns^ht^w$.
In particular,
\begin{align*}
 \text{horizontal step at height }j &: (s+t)q^{j+1},\\
 \text{up-step }j\longrightarrow j+1 &: sq^{j+2},\\
 \text{down-step }j+1\longrightarrow j &: tq^{j+1}.
\end{align*}
The two horizontal colors are distinct choices even when $s=t=1$.

\begin{example}
Take $\lambda=(5,4,2)$ and $\mu=(2,1,0)$. The row lengths are $(3,3,2)$,
the overlaps are $(2,1)$, and $(n,h,w)=(8,3,5)$. Its boundary words are
\[
 U=\mathrm{NENENEEE},\qquad L=\mathrm{EENEENEN}.
\]
The interior gaps are $(1,1,1,1,2,1,1)$, whose sum is $8$. After removing
the forced endpoint steps, the Motzkin heights are
$(0,0,0,0,1,0,0)$. The weights give $q^8s^3t^5$, as required.
\end{example}

\subsection{First-return decomposition and the Jacobi fraction}
Let $E_j(q;s,t)$ count excursions beginning and ending at height $j$ and
never going below $j$, with the weights just described. Their absolute
height matters because the power of $q$ depends on height.

A nonempty excursion begins either with a horizontal step followed by
another excursion at height $j$, or with an up-step, an excursion at height
$j+1$, a down-step, and another excursion at height $j$. Therefore
\[
 E_j=1+(s+t)q^{j+1}E_j+stq^{2j+3}E_{j+1}E_j.
\]
All nonconstant step weights have positive $q$-degree, so the equation is
valid in $\Z[s,t][[q]]$ and yields
\begin{equation}
 E_j=\frac{1}{1-(s+t)q^{j+1}-stq^{2j+3}E_{j+1}}.
 \label{eq:Erecurrence}
\end{equation}

\begin{theorem}[The connected four-variable building block]
\label{thm:Jfraction}
The connected-diagram generating function is
\begin{equation}
 P(q;s,t)=\cfrac{stq}{1-(s+t)q-
 \cfrac{stq^3}{1-(s+t)q^2-
 \cfrac{stq^5}{1-(s+t)q^3-
 \cfrac{stq^7}{1-(s+t)q^4-\ddots}}}}.
 \label{eq:Jfraction}
\end{equation}
Together with \eqref{eq:sequence}, this counts all normalized diagrams by
area, rows, columns, and connected components.
\end{theorem}
\begin{proof}
The bijection gives $P(q;s,t)=stqE_0$. Iterating
\eqref{eq:Erecurrence} proves \eqref{eq:Jfraction}. For any fixed power
$q^n$, only finitely many diagrams and path heights are possible: every
step has positive area weight. Hence the successive finite fractions
stabilize coefficientwise and define the stated formal series.
\end{proof}

This derivation includes the area statistic explicitly. Using an unweighted
Motzkin enumeration here would count a different quantity, essentially
semiperimeter, and would not establish the selected formula.

\section{Contraction to the conjectured repeated-power fraction}
\label{sec:contraction}

For the area specialization $s=t=1$, define the paired tails
\begin{equation}
 S_j=\frac{1}{1-\dfrac{q^j}{1-q^jS_{j+1}}}\qquad(j\ge1).
 \label{eq:Srecurrence}
\end{equation}
The coefficientwise interpretation is obtained by setting a sufficiently
deep tail to $1$ and then taking the depth to infinity. Each tail is a
unit with constant term $1$. Section~\ref{sec:truncation} will give an exact
bound on the resulting stabilization.

\begin{lemma}[Elementary contraction]
\label{lem:contraction}
If $R_j=S_j-1$, then
\begin{equation}
 R_j=\frac{q^j}{1-2q^j-q^jR_{j+1}}=q^jE_{j-1}(q;1,1).
 \label{eq:contraction}
\end{equation}
\end{lemma}
\begin{proof}
Combining the two fractions in \eqref{eq:Srecurrence} gives
\[
 S_j=\frac{1-q^jS_{j+1}}{1-q^j-q^jS_{j+1}}.
\]
Subtracting $1$ and substituting $S_{j+1}=1+R_{j+1}$ gives the first
identity in \eqref{eq:contraction}. On putting
$R_{j+1}=q^{j+1}E_j$, the denominator becomes
$1-2q^j-q^{2j+1}E_j$, exactly the denominator of
$E_{j-1}$ in \eqref{eq:Erecurrence}. The finite-depth identities, followed
by coefficientwise limits, prove the second equality.
\end{proof}

\begin{proof}[Proof of Theorem~\ref{thm:main}]
The expanded tail $S_1$ is $C(q)$ from \eqref{eq:mainC}. By
Lemma~\ref{lem:contraction} and Theorem~\ref{thm:Jfraction},
\[
 C(q)-1=qE_0(q;1,1)=P(q).
\]
Proposition~\ref{prop:sequence}, specialized to $s=t=u=1$, now gives
\[
 A(q)=\frac{1}{1-P(q)}=\frac{1}{2-C(q)}.
\]
Every denominator used has constant term $1$. All operations are therefore
valid as formal power-series operations, independently of numerical or
analytic convergence.
\end{proof}

\begin{corollary}[Refined repeated-power fraction]
\label{cor:refinedS}
Define
\begin{equation}
 S_j(q;s,t)=\frac{1}{1-\dfrac{sq^j}{1-tq^jS_{j+1}(q;s,t)}}.
 \label{eq:refinedS}
\end{equation}
Then
\begin{equation}
 P(q;s,t)=t\bigl(S_1(q;s,t)-1\bigr),\qquad
 \Aall(q;s,t,u)=\frac{1}{1-ut\bigl(S_1(q;s,t)-1\bigr)}.
 \label{eq:refinedSresult}
\end{equation}
\end{corollary}
\begin{proof}
Here $R_j=S_j-1$ satisfies
\[
 R_j=\frac{sq^j}{1-(s+t)q^j-tq^jR_{j+1}}
     =sq^jE_{j-1}(q;s,t).
\]
Thus $tR_1=stqE_0=P$. Apply \eqref{eq:sequence}.
\end{proof}

The symmetry in rows and columns is manifest in the Jacobi form
\eqref{eq:Jfraction}, even though it is not manifest in
\eqref{eq:refinedS}. It also follows geometrically by transposing the
diagram.

\section{A closed formula for fixed rows, columns, and components}
\label{sec:dimensions}

The area can be summed out while retaining the dimensions. For each fixed
$h,w$, all diagrams lie in an $h$-by-$w$ rectangle, so their areas are at
most $hw$. Thus setting $q=1$ is meaningful coefficientwise as a formal
series in $s,t$, even though it is not a substitution into a convergent
area generating function at $q=1$.

At $q=1$ the excursion weights no longer depend on height. All $E_j$ become
the same series $E$, with
\[
 E=1+(s+t)E+stE^2.
\]
Since $P(1;s,t)=stE$, write $P_*(s,t)=P(1;s,t)$ to obtain
\begin{equation}
 P_* =\frac{st}{1-s-t-P_*}
     =\frac{1-s-t-\sqrt{(1-s-t)^2-4st}}{2},
 \label{eq:dimensionGF}
\end{equation}
where the square root has constant term $1$. Consequently
\[
 \Aall(1;s,t,u)=\frac{1}{1-uP_*(s,t)}.
\]

\begin{theorem}[Fixed-dimension general term]
\label{thm:dimensionformula}
Let $d_{h,w,k}$ count normalized skew diagrams with exactly $h$ rows, $w$
columns, and $k$ edge-connected components, without fixing their area. For
$h,w\ge k\ge1$,
\begin{equation}
 \boxed{\displaystyle
 d_{h,w,k}=\frac{k}{h+w-k}
 \binom{h+w-k}{h}\binom{h+w-k}{w}.}
 \label{eq:dimensionformula}
\end{equation}
It is zero if $k\ge1$ and either $h<k$ or $w<k$. The empty case is
$d_{0,0,0}=1$, and no nonempty diagram has $k=0$.
\end{theorem}
\begin{proof}
By unique concatenation, $d_{h,w,k}=[s^ht^w]P_*^k$. Introduce a temporary
variable $z$ and the unique series
\[
 Y=\frac{z}{1-s-t-Y}.
\]
Lagrange coefficient extraction gives, for $m\ge k$,
\begin{align}
 [z^m]Y^k
 &=\frac{k}{m}[y^{m-k}](1-s-t-y)^{-m}\notag\\
 &=\frac{k}{m}\binom{2m-k-1}{m-k}(1-s-t)^{-(2m-k)}.
 \label{eq:lagrange}
\end{align}
Substitute $z=st$, so $Y=P_*$. Extracting the remaining powers of $s$ and
$t$ gives
\begin{equation}
 d_{h,w,k}=k(h+w-k-1)!
 \sum_{m=k}^{\min(h,w)}
 \frac{1}{m!(m-k)!(h-m)!(w-m)!}.
 \label{eq:dimensionintermediate}
\end{equation}
The sum equals
\begin{align*}
 \frac{1}{h!(w-k)!}
 \sum_m\binom{h}{m}\binom{w-k}{m-k}
 &=\frac{1}{h!(w-k)!}\binom{h+w-k}{w}\\
 &=\frac{(h+w-k)!}{h!w!(h-k)!(w-k)!},
\end{align*}
by Vandermonde's identity. Substitution in
\eqref{eq:dimensionintermediate} is exactly
\eqref{eq:dimensionformula}. The zero cases follow because each nonempty
component uses at least one row and one column.
\end{proof}

For $k=1$, this recovers the usual fixed-dimension parallelogram-polyomino
count. For $h=w=2$, there are three connected diagrams and one diagram
with two components. They do not all have the same area, which is why this
formula should not be confused with the area coefficients $a_n$.

\section{An independent solution by a ratio of \texorpdfstring{$q$}{q}-series}
\label{sec:qseries}

The boundary-path proof already proves the selected identity. A second
algebraic route is valuable both as a check and as an alternative formula
for exact coefficient computation. Define
\begin{equation}
 \Phi(z;q)=\sum_{r\ge0}
 \frac{(-1)^r z^r q^{r(r+1)/2}}{(q;q)_r^2},
 \qquad (q;q)_r=\prod_{j=1}^r(1-q^j),\quad(q;q)_0=1.
 \label{eq:Phi}
\end{equation}
This is a formal power series in $q$ with polynomial coefficients in $z$.
It is also analytic for $|q|<1$ and fixed $z$: on compact subdisks the
finite products in the denominator are bounded away from zero, whereas
the numerator decays quadratically in the exponent of $r$.

The connected ratio below is consistent with Bala's treatment and the
Delest--F\'edou enumeration \cite{bala,delestfedou}. We derive it directly
from the overlap coordinates, so no continued-fraction identity from those
sources is required for this proof.

\begin{theorem}[The row-refined ratio]
\label{thm:qseries}
One has
\begin{equation}
 1+P(q;s,1)=\frac{\Phi(sq;q)}{\Phi(s;q)},
 \label{eq:Pratio}
\end{equation}
and hence
\begin{equation}
 \Aall(q;s,1,u)=
 \frac{\Phi(s;q)}{(1+u)\Phi(s;q)-u\Phi(sq;q)}.
 \label{eq:AratioRefined}
\end{equation}
In particular, if $D(q)=\Phi(1;q)$ and $N(q)=\Phi(q;q)$, then
\begin{equation}
 \boxed{\displaystyle A(q)=\frac{D(q)}{2D(q)-N(q)}.}
 \label{eq:Aratio}
\end{equation}
\end{theorem}
\begin{proof}
For $i\ge1$, let $X_i$ count all possible continuations below a
previously given row of length $i$, not counting that row's own weight.
The empty continuation is allowed. Appending a row of length $j$ has
weight $sq^j$, and there are $\min(i,j)$ positive overlap choices.
Therefore
\begin{equation}
 X_i=1+s\sum_{j\ge1}\min(i,j)q^jX_j.
 \label{eq:transfer}
\end{equation}
There is a unique family of formal solutions with constant terms $1$:
every appended row increases area, so iteration determines one more
coefficient at a time. Choosing the initial row separately gives
\begin{equation}
 P(q;s,1)=s\sum_{j\ge1}q^jX_j=X_1-1.
 \label{eq:firstrow}
\end{equation}

Set $Z_i=\Phi(sq^i;q)$ for $i\ge0$. If
$c_r=(-1)^rq^{r(r+1)/2}/(q;q)_r^2$, then
$c_r(1-q^r)^2=-q^rc_{r-1}$ for $r\ge1$. Comparing coefficients of $s^r$
now gives
\begin{equation}
 Z_{i+1}-(2-sq^i)Z_i+Z_{i-1}=0\qquad(i\ge1).
 \label{eq:Zrecurrence}
\end{equation}
Let $\Delta_i=Z_i-Z_{i-1}$. Both $Z_i\to1$ and $\Delta_i\to0$
coefficientwise as $i\to\infty$. Equation~\eqref{eq:Zrecurrence} says
\[
 \Delta_i-\Delta_{i+1}=sq^iZ_i.
\]
Summing first over $i\ge a$, and then over $1\le a\le k$, gives
\begin{align*}
 \Delta_a&=s\sum_{j\ge a}q^jZ_j,\\
 Z_k-Z_0&=s\sum_{j\ge1}\min(k,j)q^jZ_j.
\end{align*}
Every sum is coefficientwise finite. The series $Z_0=\Phi(s;q)$ has
constant term $1$, so division by it is valid. The family $Z_k/Z_0$ solves
\eqref{eq:transfer}, and uniqueness implies $X_k=Z_k/Z_0$.
Equation~\eqref{eq:firstrow} proves \eqref{eq:Pratio}. Finally apply the
component identity \eqref{eq:sequence} and simplify.
\end{proof}

Expanding the two series in \eqref{eq:Aratio} explicitly gives
\[
 D(q)=\sum_{r\ge0}\frac{(-1)^rq^{r(r+1)/2}}{(q;q)_r^2},
 \qquad
 N(q)=\sum_{r\ge0}\frac{(-1)^rq^{r(r+3)/2}}{(q;q)_r^2}.
\]
Thus the selected continued fraction and a ratio of explicit infinite sums
have now been derived from two different encodings of the same diagrams.

\subsection{A finite transfer matrix and a different continued fraction}
This construction will also supply an exact numerical certificate. Restrict
every row length to be at most $H$, and let $P_H(q)$ count the connected
diagrams satisfying that restriction. Set
\[
 M_H=(\min(i,j))_{1\le i,j\le H},\qquad
 D_H=\operatorname{diag}(q,q^2,\ldots,q^H),
\]
and let $\mathbf e$ be the all-ones column vector. The transfer equation
becomes
\begin{equation}
 \mathbf X=\mathbf e+M_HD_H\mathbf X,
 \qquad P_H=\mathbf e^{\mathsf T}D_H\mathbf X.
 \label{eq:finiteTransfer}
\end{equation}
The matrix $M_H$ is $LL^{\mathsf T}$, where $L$ is lower triangular with
all entries on and below its diagonal equal to $1$. Its inverse is
\begin{equation}
 B_H=M_H^{-1}=
 \begin{pmatrix}
 2&-1& & &\\
 -1&2&-1& &\\
 &\ddots&\ddots&\ddots&\\
 &&-1&2&-1\\
 &&&-1&1
 \end{pmatrix}
 \quad(H\ge2),\qquad B_1=(1).
 \label{eq:tridiagonal}
\end{equation}
Since $B_H\mathbf e$ is the first standard basis vector, and
$X_1=1+P_H$, it follows that
\begin{equation}
 1+P_H(q)=(B_H-D_H)^{-1}_{11}.
 \label{eq:inverseentry}
\end{equation}
Successive Schur complements give the entirely finite fraction
\begin{equation}
 1+P_H(q)=\cfrac{1}{2-q-\cfrac{1}{2-q^2-\cfrac{1}{\ddots-
 \cfrac{1}{2-q^{H-1}-\cfrac{1}{1-q^H}}}}}
 \qquad(H\ge2),
 \label{eq:boundedCF}
\end{equation}
with $1+P_1(q)=1/(1-q)$. Notice the final denominator $1-q^H$, not
$2-q^H$. This boundary condition matters in both exact computation and
error bounds.

\section{Sharp truncation and an exact coefficient algorithm}
\label{sec:truncation}

\subsection{The first omitted term is a square}
Let $C^{(m)}(q)$ be the paired fraction obtained by setting
$S_{m+1}^{(m)}=1$ and applying \eqref{eq:Srecurrence} for
$j=m,m-1,\ldots,1$. Put
\[
 A^{(m)}(q)=\frac{1}{2-C^{(m)}(q)}.
\]
For $m=0$ this means $C^{(0)}=A^{(0)}=1$.

\begin{theorem}[Exact truncation order]
\label{thm:truncation}
For every integer $m\ge0$,
\begin{align}
 C(q)-C^{(m)}(q)&=q^{(m+1)^2}+O\bigl(q^{(m+1)^2+1}\bigr),
 \label{eq:truncC}\\
 A(q)-A^{(m)}(q)&=q^{(m+1)^2}+O\bigl(q^{(m+1)^2+1}\bigr).
 \label{eq:truncA}
\end{align}
Here $O(q^r)$ means a formal series divisible by $q^r$.
Consequently $m=\lfloor\sqrt N\rfloor$ pairs suffice to compute every
coefficient through $q^N$.
\end{theorem}
\begin{proof}
Write the paired-tail map as
\[
 \phi_j(z)=\frac{1-q^jz}{1-q^j-q^jz}.
\]
For any two formal units $a,b$ with constant term $1$, direct subtraction
gives
\begin{equation}
 \phi_j(a)-\phi_j(b)=
 \frac{q^{2j}(a-b)}
 {(1-q^j-q^ja)(1-q^j-q^jb)}.
 \label{eq:phidifference}
\end{equation}
At the first omitted tail,
\[
 S_{m+1}-1=q^{m+1}+O(q^{m+2}).
\]
Applying \eqref{eq:phidifference} for $j=m,m-1,\ldots,1$ increases the
valuation by
$2\sum_{j=1}^m j=m(m+1)$ and leaves the leading coefficient equal to $1$.
This proves \eqref{eq:truncC}, since
$m(m+1)+(m+1)=(m+1)^2$. It also proves directly that the finite fractions
stabilize coefficientwise.

Finally,
\[
 \frac{1}{2-C}-\frac{1}{2-C^{(m)}}=
 \frac{C-C^{(m)}}{(2-C)(2-C^{(m)})}.
\]
Both denominator factors have constant term $1$, so the same valuation and
leading coefficient hold for $A-A^{(m)}$. Since
$(\lfloor\sqrt N\rfloor+1)^2>N$, the stated depth is sufficient.
\end{proof}

There is a geometric explanation. Contracting a fraction of depth $m$
allows Motzkin heights at most $m-1$, or boundary gaps at most $m$. The
least-area omitted excursion rises to height $m$ and immediately descends,
with no horizontal steps. Its diagram is the $(m+1)$-by-$(m+1)$ square,
with area $(m+1)^2$. It is unique at that least area, explaining the leading
coefficient $1$. Additional components would increase the area, so the
same first discrepancy holds for all diagrams.

For $N=1000$, only $31$ pairs, that is, $62$ individual partial numerators,
are needed. The first missing term has degree $32^2=1024$. This statement
is exact, not an empirical rule about rapid numerical convergence.

\subsection{Continuants avoid repeated series inversion}
Write a finite tail as $U/V$, with $U(0)=V(0)=1$. Then
\[
 \phi_j(U/V)=\frac{V-q^jU}{V-q^jV-q^jU}.
\]
Thus the finite polynomials $U,V$ for $C^{(m)}=U/V$ can be computed by
starting with $U=V=1$ and performing the simultaneous updates
\begin{equation}
 \begin{split}
 U_{\rm new}&=V-q^jU,\\
 V_{\rm new}&=V-q^jV-q^jU,
 \end{split}
 \qquad j=m,m-1,\ldots,1.
 \label{eq:continuants}
\end{equation}
All polynomials may be truncated after degree $N$. At the end,
\begin{equation}
 A^{(m)}=\frac{V}{2V-U}.
 \label{eq:continuantA}
\end{equation}
Only one series division is needed to obtain $A$.

If $B(q)=\sum b_jq^j=2V-U$ and $V(q)=\sum v_jq^j$, then $b_0=1$ and
\begin{equation}
 a_0=1,\qquad
 a_n=v_n-\sum_{j=1}^{n}b_ja_{n-j}\quad(1\le n\le N).
 \label{eq:divisionrecurrence}
\end{equation}
The shift-and-subtract updates cost $O(N\sqrt N)$ integer arithmetic
operations. Ordinary dense division costs $O(N^2)$ such operations. The
continuant method uses $O(N)$ integer storage. These are arithmetic-operation
bounds, not bit-complexity bounds; the integers grow with $N$.

The included implementation also evaluates the overlap sums by a different
dynamic program. Let $f_{n,i}$ count diagrams of area $n$ whose last row has
length $i$. Initialize $f_{i,i}=1$, and append a row of length $j$ with
multiplier $\min(i,j)+\varepsilon$, where $\varepsilon=0$ for connected
shapes and $\varepsilon=1$ for all shapes. The transition sum can be written
\begin{equation}
 \sum_i f_{n,i}(\min(i,j)+\varepsilon)
 =\sum_{i\le j}i f_{n,i}
  +j\sum_{i>j}f_{n,i}
  +\varepsilon\sum_i f_{n,i}.
 \label{eq:prefixDP}
\end{equation}
Prefix sums give an $O(N^2)$-operation dynamic program, using $O(N^2)$
storage in the supplied straightforward implementation. It never evaluates
a continued fraction.

\section{Rigorous exponential asymptotics}
\label{sec:asymptotics}

\subsection{Analytic control from the row-transfer kernel}
The formal proofs imply exact enumeration. To derive asymptotics, we now
establish an analytic disk large enough to contain the dominant pole.
For $s=1$, regard the transfer kernel
\[
 K(q)_{ij}=\min(i,j)q^j\qquad(i,j\ge1)
\]
as an operator on the Banach space of bounded sequences, with the supremum
norm. For $|q|\le r<1$,
\begin{equation}
 \opnorm{K(q)}\le\sum_{j\ge1}jr^j=\frac{r}{(1-r)^2}=:\eta(r).
 \label{eq:kernelbound}
\end{equation}
The operator-valued series defining $K$ is analytic in that disk. In
particular, if
\begin{equation}
 r<r_0:=\frac{3-\sqrt5}{2}=0.381966\ldots,
 \label{eq:rzero}
\end{equation}
then $\eta(r)<1$, so
\[
 \mathbf X=(I-K)^{-1}\mathbf e=\sum_{m\ge0}K^m\mathbf e
\]
converges uniformly in operator norm. The functional
$\mathbf f(q)^{\mathsf T}\mathbf X=\sum_{j\ge1}q^jX_j$ has norm bounded
by $r/(1-r)$. Hence $P(q)$ is analytic for $|q|<r_0$.

\begin{lemma}[Existence of a positive singularity before $1/3$]
\label{lem:rho}
There is a unique $\rho\in(0,1/3)$ such that $P(\rho)=1$, and
$P'(\rho)>0$.
\end{lemma}
\begin{proof}
The coefficients of $P$ are nonnegative, and $p_1=1$. Thus $P$ is strictly
increasing on $(0,r_0)$, with $P(0)=0$ and positive derivative. Its first
seven terms are
\[
 q+2q^2+4q^3+9q^4+20q^5+46q^6+105q^7.
\]
At $q=1/3$ their sum is $245/243>1$. Analyticity at $1/3$, continuity,
and strict monotonicity now prove the assertion.
\end{proof}

\begin{theorem}[The asymptotic formula]
\label{thm:asymptotic}
Let $\rho$ be as in Lemma~\ref{lem:rho}. There is an $R>\rho$ such that
\begin{equation}
 a_n=\kappa\rho^{-n}+O(R^{-n}),\qquad
 \kappa=\frac{1}{\rho P'(\rho)}>0.
 \label{eq:asymptotic}
\end{equation}
In particular, $\lim_{n\to\infty}a_n^{1/n}=1/\rho$, and
$a_{n+1}/a_n\to1/\rho$.
\end{theorem}
\begin{proof}
For $|q|<\rho$,
$|P(q)|\le P(|q|)<1$, so $A=1/(1-P)$ has no pole. If $|q|=\rho$ and
$P(q)=1$, equality holds in the triangle inequality for
$\sum_{n\ge1}p_nq^n$. The terms of degrees $1$ and $2$ have positive
coefficients; their phases must agree. Writing $q=\rho e^{i\theta}$ gives
$e^{i\theta}=e^{2i\theta}$, hence $q=\rho$. Thus $\rho$ is the only zero
of $1-P$ on its circle.

At $\rho$, the derivative $P'(\rho)$ is positive, so the pole is simple.
Because $P$ is analytic in a larger disk and analytic zeros are isolated,
there are numbers $\rho<R<R'<r_0$ for which $1-P$ has no other zero in
$|q|\le R'$. Its principal part at $\rho$ is
\[
 \frac{1}{\rho P'(\rho)}\frac{1}{1-q/\rho}.
\]
Subtracting this function from $A$ removes the pole. The difference is
analytic on $|q|<R'$, so Cauchy's coefficient bound on $|q|=R$ is
$O(R^{-n})$. This proves \eqref{eq:asymptotic}. The two stated limits
follow because $\kappa>0$.
\end{proof}

For reference, $D(q)=\Phi(1;q)$ does not vanish in $|q|<r_0$. Indeed, for
fixed $q$ there the bounded sequence $Z_i=\Phi(q^i;q)$ satisfies
$\mathbf Z=D\mathbf e+K\mathbf Z$, by the proof of
Theorem~\ref{thm:qseries}. If $D=0$, invertibility of $I-K$ forces
$\mathbf Z=0$, contradicting $Z_i\to1$. Therefore the equation for $\rho$
can equivalently be written
\begin{equation}
 2\Phi(1;\rho)-\Phi(\rho;\rho)=0,
 \label{eq:rootqseries}
\end{equation}
with a nonzero numerator in \eqref{eq:Aratio}.

\subsection{A rational certificate for the dominant root}
High-precision evaluation suggests
\begin{align}
 \rho&\approx0.3195967180593874655186028919827139236,\label{eq:rhonum}\\
 \rho^{-1}&\approx3.128943269730886252277447995387754161,\label{eq:growthnum}\\
 \kappa&\approx0.2893214639716513230848453480726559061.
 \label{eq:kappanum}
\end{align}
These displayed extended decimals are numerical estimates. The following
strict interval is separately certified using exact rational arithmetic:
\begin{equation}
 \boxed{\quad
 0.31959671805938746551<\rho<0.31959671805938746552.
 \quad}
 \label{eq:rootinterval}
\end{equation}

Here is the complete bound behind that certificate. For real
$0<q<r_0$, define
\begin{align}
 \tau_H(q)&=\sum_{j>H}q^j=\frac{q^{H+1}}{1-q},\label{eq:tau}\\
 \delta_H(q)&=\sum_{j>H}jq^j
 =\frac{q^{H+1}\bigl((H+1)-Hq\bigr)}{(1-q)^2}.
 \label{eq:delta}
\end{align}

\begin{proposition}[An explicit bounded-row error]
\label{prop:rowerror}
For $0<q<r_0$,
\begin{equation}
 0\le P(q)-P_H(q)\le
 \frac{\tau_H(q)}{1-\eta(q)}+
 \frac{q}{1-q}\frac{\delta_H(q)}{(1-\eta(q))^2}
 =:\mathcal E_H(q).
 \label{eq:rowerror}
\end{equation}
\end{proposition}
\begin{proof}
Extend the finite kernel to bounded infinite sequences by suppressing all
columns with index $j>H$; denote this operator by $K_H$. Its first $H$
coordinates satisfy the finite system \eqref{eq:finiteTransfer}. The
operator norms obey
\[
 \opnorm{K_H}\le\eta(q),\qquad
 \opnorm{K-K_H}\le\delta_H(q).
\]
Let $\mathbf f$ be the functional with entries $q^j$, and $\mathbf f_H$
its truncation. Then
\[
 \opnorm{\mathbf f-\mathbf f_H}=\tau_H(q),\qquad
 \opnorm{\mathbf f_H}\le\frac{q}{1-q}.
\]
Using the resolvent identity,
\[
 (I-K)^{-1}-(I-K_H)^{-1}
 =(I-K)^{-1}(K-K_H)(I-K_H)^{-1},
\]
and $\opnorm{\mathbf e}=1$, we obtain the upper bound in
\eqref{eq:rowerror}. The lower bound also follows directly because
$P_H$ counts a subset of the diagrams counted by $P$.
\end{proof}

Let $q_-$ and $q_+$ denote the two rational endpoints in
\eqref{eq:rootinterval}. The script \texttt{verify.py} evaluates
\eqref{eq:boundedCF} with $H=55$ using Python's exact rational type and
checks
\begin{equation}
 P_{55}(q_-)+\mathcal E_{55}(q_-)<1,
 \qquad P_{55}(q_+)>1.
 \label{eq:certificatechecks}
\end{equation}
The error allowance at the lower endpoint is about
$7.446\times10^{-26}$. No floating-point decision is used in
\eqref{eq:certificatechecks}. Proposition~\ref{prop:rowerror}, combined
with monotonicity of $P$, proves \eqref{eq:rootinterval}. The interval has
width $10^{-20}$; this statement should not be confused with certification
of every additional decimal printed in \eqref{eq:rhonum}--\eqref{eq:kappanum}.

\section{The number of connected components in a large diagram}
\label{sec:componentsCLT}

The component marker gives more than a counting refinement. Choose a
diagram uniformly among the $a_n$ diagrams of area $n$, for $n\ge1$, and
let $K_n$ be its number of edge-connected components. By
\eqref{eq:sequence},
\begin{equation}
 \E[u^{K_n}]=\frac{[q^n](1-uP(q))^{-1}}{a_n}.
 \label{eq:componentPGF}
\end{equation}

\begin{theorem}[Linear mean, linear variance, and a Gaussian limit]
\label{thm:CLT}
Define
\begin{equation}
 \mu=\frac{1}{\rho P'(\rho)},\qquad
 \sigma^2=\mu^2-\mu+\frac{P''(\rho)}{\rho P'(\rho)^3}.
 \label{eq:CLTconstants}
\end{equation}
Then $\sigma^2>0$,
\begin{equation}
 \E[K_n]=\mu n+O(1),\qquad
 \Var(K_n)=\sigma^2n+O(1),
 \label{eq:CLTmoments}
\end{equation}
and
\begin{equation}
 \frac{K_n-\mu n}{\sigma\sqrt n}
 \ \xrightarrow{\ d\ }\ \mathcal N(0,1).
 \label{eq:CLT}
\end{equation}
Numerically, $\mu=\kappa\approx0.2893214639716513231$ and
$\sigma^2\approx0.2382004452790812020$.
\end{theorem}
\begin{proof}
The implicit function theorem gives an analytic function $\rho(u)$ near
$u=1$ such that
\[
 uP(\rho(u))=1,\qquad \rho(1)=\rho.
\]
Choose a circle $|q|=R>\rho$ as in Theorem~\ref{thm:asymptotic}. For
complex $u$ in a sufficiently small neighborhood of $1$, continuity and
the argument principle show that $1-uP(q)$ has just the one simple zero
$\rho(u)$ inside that circle and no zero on it. Subtracting the moving
principal part and applying Cauchy's bound uniformly gives
\begin{equation}
 [q^n]\frac{1}{1-uP(q)}
 =\frac{\rho(u)^{-n}}{u\rho(u)P'(\rho(u))}+O(R^{-n}).
 \label{eq:uniformpole}
\end{equation}
The leading factor is analytic and nonzero near $u=1$.

Set $r(t)=\rho(e^t)$ and $L(t)=\log(\rho/r(t))$, taking local analytic
branches. Equation~\eqref{eq:componentPGF} and the uniform estimate imply
\begin{equation}
 \log\E[e^{tK_n}]=nL(t)+G(t)+O(\theta^n),\qquad 0<\theta<1,
 \label{eq:quasipowers}
\end{equation}
locally uniformly in complex $t$, for an analytic $G$ independent of $n$.
Differentiate $P(r(t))=e^{-t}$ at zero to obtain
\[
 r'(0)=-\frac{1}{P'(\rho)},\qquad
 r''(0)=\frac{1}{P'(\rho)}-
        \frac{P''(\rho)}{P'(\rho)^3}.
\]
It follows that $L'(0)=\mu$ and $L''(0)=\sigma^2$, proving
\eqref{eq:CLTmoments} by differentiation of the logarithmic moment
generating function.

To verify strict positivity, set $w_j=p_j\rho^j$. Since $P(\rho)=1$, these
numbers form a probability distribution on positive integers. Its mean
$m$ and variance $v$ satisfy
\[
 m=\rho P'(\rho),\qquad
 v=\rho^2P''(\rho)+\rho P'(\rho)-\bigl(\rho P'(\rho)\bigr)^2.
\]
A direct simplification gives $\sigma^2=v/m^3$. Both $w_1$ and $w_2$ are
positive, so this distribution is not concentrated at one point; hence
$v>0$.

Finally, substitute $t=iz/(\sigma\sqrt n)$ in
\eqref{eq:quasipowers} for fixed real $z$ and subtract the centering term
$iz\mu\sqrt n/\sigma$. Taylor expansion gives a limit of $-z^2/2$ for the
logarithm of the characteristic function. Its exponential therefore tends
to $e^{-z^2/2}$, proving \eqref{eq:CLT} by the continuity theorem for
characteristic functions.
\end{proof}

The equality between the mean-component density $\mu$ and the leading
coefficient $\kappa$ in the area asymptotic is a consequence of the same
simple-pole calculation, not a numerical coincidence. These consequences
are derived here without a claim that their general analytic mechanism is
new.

\section{Exact checks, numerical checks, and reproducibility}
\label{sec:verification}

\subsection{Initial values and component polynomials}
Table~\ref{tab:initial} compares the two area sequences obtained from the
proof. The connected sequence is A006958, with its constant term set to
zero; the all-diagram sequence is A225114, including the empty object
\cite{oeis006958,oeis225114}.

\begin{table}[htbp]
\centering
\begin{tabular}{@{}r r r@{}}
\toprule
Area $n$ & Connected diagrams $p_n$ & All diagrams $a_n$\\
\midrule
0 & 0 & 1\\
1 & 1 & 1\\
2 & 2 & 3\\
3 & 4 & 9\\
4 & 9 & 28\\
5 & 20 & 87\\
6 & 46 & 272\\
7 & 105 & 850\\
8 & 242 & 2659\\
9 & 557 & 8318\\
10 & 1285 & 26025\\
11 & 2964 & 81427\\
12 & 6842 & 254777\\
\bottomrule
\end{tabular}
\caption{Initial exact coefficients.}\label{tab:initial}
\end{table}

For example, direct overlap enumeration gives
\begin{align*}
 a_1(u)&=u,\\
 a_2(u)&=2u+u^2,\\
 a_3(u)&=4u+4u^2+u^3,\\
 a_4(u)&=9u+12u^2+6u^3+u^4,\\
 a_5(u)&=20u+34u^2+24u^3+8u^4+u^5,\\
 a_6(u)&=46u+92u^2+83u^3+40u^4+10u^5+u^6.
\end{align*}
Their coefficients agree with $[q^n]P(q)^k$ for each component count $k$.
Setting $u=1$ recovers the corresponding values in the table.

\subsection{The exact verification suite}
The following checks were executed successfully by
\texttt{python3 verify.py --limit 1000}. The summary is preserved in
\texttt{verification\_results.json}.
\begin{enumerate}[leftmargin=1.8em,itemsep=3pt]
\item The recorded candidate-list digest and random draw were replayed
without creating a new draw.
\item All $31$ initial terms displayed in A225114 and all $32$ positive-area
terms displayed in A006958 agreed with the computed values.
\item Independent row-overlap dynamic programs, the paired continued
fraction, the explicit $q$-series ratio, and the ordered-component
convolution agreed through area $1000$.
\item The exact first omitted coefficient in
Theorem~\ref{thm:truncation} was checked at every depth from $0$ to $12$.
\item Full component polynomials were independently enumerated from
compositions through area $13$.
\item The four-variable refinement was checked through area $20$ at all
$27$ triples $(s,t,u)\in\{0,1,2\}^3$, against direct weighted row placements.
These are specialized tests, not a claim of symbolic verification at all
four-variable coefficients through area $1000$.
\item Direct enumeration of the two boundary paths agreed with the
connected area counts through area $10$. Since
$n\ge h+w-1$ by \eqref{eq:areasum}, enumerating boundary lengths through
$11$ covers every diagram in that range.
\item The dimension formula \eqref{eq:dimensionformula} agreed with
independent boundary-path enumeration and concatenation for every
$h,w\le5$ and $1\le k\le5$.
\item The dominant-root inequalities \eqref{eq:certificatechecks} were
verified with exact rational arithmetic, and the seven-term lower bound
$P(1/3)>245/243$ was checked.
\end{enumerate}

The program reported \texttt{ALL EXACT CHECKS PASSED}. These finite checks
are designed to expose indexing, weighting, truncation, and implementation
errors. They are not used as a replacement for any all-orders proof in
Sections~\ref{sec:rows}--\ref{sec:componentsCLT}.

\subsection{Numerical diagnostics for the asymptotics}
The optional script \texttt{asymptotics.py} uses \texttt{mpmath} at high
precision to estimate $\rho$, $P'(\rho)$, and $P''(\rho)$. It compares
continued-fraction depths $40$ and $60$ as a stability check. Agreement of
those evaluations is a numerical diagnostic, not the rational certificate;
only \eqref{eq:rootinterval} is certified in that way.

The moments can also be recovered from integer coefficient arrays, using
\begin{equation}
 \E[K_n]=\frac{[q^n](A^2-A)}{a_n},\qquad
 \E[K_n(K_n-1)]=\frac{2[q^n](A^3-2A^2+A)}{a_n}.
 \label{eq:exactmoments}
\end{equation}
These follow by differentiating $(1-uP)^{-1}$ and using $P=1-1/A$.
Table~\ref{tab:asymptotic} compares the resulting statistics with their
limits.

\begin{table}[htbp]
\centering\small
\begin{tabular}{@{}r c c c@{}}
\toprule
$n$ & $a_n\rho^n$ & $\E[K_n]/n$ & $\Var(K_n)/n$\\
\midrule
10 & 0.289343577662 & 0.371627281460 & 0.209156420607\\
20 & 0.289321476436 & 0.330486811475 & 0.223648203436\\
50 & 0.289321463972 & 0.305787608032 & 0.232379523182\\
100 & 0.289321463972 & 0.297554536002 & 0.235289984230\\
250 & 0.289321463972 & 0.292614692784 & 0.237036260860\\
500 & 0.289321463972 & 0.290968078378 & 0.237618353069\\
1000 & 0.289321463972 & 0.290144771175 & 0.237909399174\\
\midrule
Limit & 0.289321463972 & 0.289321463972 & 0.238200445279\\
\bottomrule
\end{tabular}
\caption{Numerical diagnostics from exact area and component counts, with
high-precision estimates of the asymptotic constants.}
\label{tab:asymptotic}
\end{table}

\subsection{Contents and scope of the accompanying archive}
The archive contains the editable \LaTeX{} source and its compiled PDF,
the frozen candidate list and selection record, the standard-library exact
verification program, the optional numerical program, and CSV/JSON outputs.
The coefficient table contains $n=0,1,\ldots,1000$. The separate component
table contains the nonzero entries for $n\le13$. The README gives build and
execution commands and distinguishes rigorous inequalities from numerical
estimates.

No proof assistant or external referee has certified this article, and no
edit has been submitted to OEIS as part of this work. The all-$n$ result
rests on the explicit mathematical arguments above; the code provides
independent reproducible checks of their consequences.

\section{Conclusion}
The generating function selected from the OEIS shortlist is correct.
Its structure has a direct combinatorial explanation: zero row overlaps
cut a diagram into an ordered sequence of connected pieces; each connected
piece is encoded by two separated monotone boundaries; their separation
is a weighted Motzkin excursion. First-return decomposition produces the
Jacobi fraction, and elementary contraction produces exactly the repeated
powers in the selected expression.

The same framework yields the full refinement
$\Aall(q;s,t,u)=(1-uP(q;s,t))^{-1}$, the dimension formula
\eqref{eq:dimensionformula}, the independent $q$-series ratio
\eqref{eq:Aratio}, the sharp truncation degree $(m+1)^2$, and the rigorous
asymptotic \eqref{eq:asymptotic}. The prior sketch is explicitly credited;
the deliverable is a complete proof and reproducible elaboration of a
formula still labeled conjectural in its OEIS entry, not a priority claim.

\appendix
\section{The unselected candidate formulas and draw record}
\label{app:candidates}

The following formulas were present in the frozen list. They are recorded
to make the selection auditable; this article makes no assertion that they
have been proved or refuted here.

For A244475, with the power-series indexing beginning at $n=3$, the
conjectured generating function is \cite{oeis244475}
\begin{equation}
 -\frac{q^3 Q(q)}{q^2+q-1},
 \qquad
 \begin{aligned}
 Q(q)={}&q^{14}+q^{13}+q^{12}+2q^{11}+3q^{10}+5q^9+8q^8\\
       &+q^7+3q^6+3q^5+2q^4+4q^3+5q^2+2q+1.
 \end{aligned}
 \label{eq:candidateStern}
\end{equation}
The OEIS attribution is Alois P.\ Heinz, 20 June 2022.

For A381190, with indexing beginning at $n=3$, the conjectured generating
function is \cite{oeis381190}
\begin{equation}
 -\frac{2q^3(4q^5+8q^4+4q^3-9q^2-8q-3)}{(q^3+q^2-1)^2},
 \label{eq:candidateDomination}
\end{equation}
attributed to Joerg Arndt, 7 January 2026.

For A391632, with indexing beginning at $n=1$, the conjectured generating
function is \cite{oeis391632}
\begin{equation}
 \frac{q(4-19q+33q^2-40q^3+8q^4)}{1-10q+29q^2-32q^3+8q^4},
 \label{eq:candidateVisibility}
\end{equation}
attributed to Andrew Howroyd, 12 January 2026.

The SHA-256 digest of the exact UTF-8 bytes of \texttt{candidates.json} is
\begin{center}\ttfamily\small
 ea6cca457acda9a1fbd92f88aa09cd9250053f0f467d75e0f7323aabc7c82d37
\end{center}
The RNG seed is given in Section~\ref{sec:selection}; the recorded output
is $0$. The exact verification program replays this draw without altering
the record. The original selection script deliberately refuses to run
when \texttt{selection.json} already exists, to avoid silently replacing
the recorded selection.

\Needspace{10\baselineskip}
\section{A minimal reproduction protocol}
\label{app:reproduce}
Compile the article twice with a standard \LaTeX{} installation:
\begin{verbatim}
pdflatex -interaction=nonstopmode -halt-on-error article.tex
pdflatex -interaction=nonstopmode -halt-on-error article.tex
\end{verbatim}
Run the exact checks with Python 3.10 or later:
\begin{verbatim}
python3 verify.py --limit 1000
\end{verbatim}
This step uses only the Python standard library. For the optional numerical
tables, install the separately listed dependency and run:
\begin{verbatim}
python3 -m pip install -r requirements.txt
python3 asymptotics.py
\end{verbatim}
The numerical script computes its own coefficient arrays through area
$1000$ using the continuant routine. Running the exact script with a
different limit overwrites its CSV coefficient output at that limit; use
$1000$ to reproduce the supplied exact table. Neither command performs a
new random selection.

\begin{thebibliography}{99}
\raggedright
\addcontentsline{toc}{section}{References}

\bibitem{oeis225114}
The OEIS Foundation Inc.,
\emph{A225114: Number of skew partitions of $n$ whose diagrams have no
empty rows and columns}.
Entry by Joerg Arndt; generating-function conjecture attributed to
Mikhail Kurkov, 3 September 2024.
\url{https://oeis.org/A225114}. Consulted 20 September 2026.

\bibitem{oeis244475}
The OEIS Foundation Inc.,
\emph{A244475: Fifth-largest distinct value in a row of Stern's diatomic
triangle}.
\url{https://oeis.org/A244475}. Consulted 20 September 2026.

\bibitem{oeis381190}
The OEIS Foundation Inc.,
\emph{A381190: Connected minimal dominating sets of $n$-trapezohedral graphs}.
\url{https://oeis.org/A381190}. Consulted 20 September 2026.

\bibitem{oeis391632}
The OEIS Foundation Inc.,
\emph{A391632: Mutual-visibility sets of $n$-Andr\'asfai graphs}.
\url{https://oeis.org/A391632}. Consulted 20 September 2026.

\bibitem{pasechnik}
Dima Pasechnik,
Answer to \emph{Generating function for A225114}, MathOverflow,
10 September 2024.
\url{https://mathoverflow.net/questions/478171/generating-function-for-a225114}.
Consulted 20 September 2026.

\bibitem{oeis006958}
The OEIS Foundation Inc.,
\emph{A006958: Number of parallelogram polyominoes with $n$ cells}.
Includes continued-fraction formulas attributed to Paul D.\ Hanna and
$q$-series formulas attributed to Peter Bala.
\url{https://oeis.org/A006958}. Consulted 20 September 2026.

\bibitem{bala}
Peter Bala,
\emph{Fountains of coins and skew Ferrers diagrams}, 26 July 2019, 6 pages.
Note linked from OEIS A161492 and A006958.
\url{https://oeis.org/A161492/a161492_1.pdf}.

\bibitem{delestfedou}
M.-P.\ Delest and J.-M.\ F\'edou,
\emph{Enumeration of skew Ferrers diagrams},
Discrete Mathematics \textbf{112} (1993), 65--79.
\url{https://doi.org/10.1016/0012-365X(93)90224-H}.
Historical background; the identities used here are independently derived
in the article.

\end{thebibliography}

\end{document}
